\documentclass[11pt,a4paper]{article}
\usepackage[T1]{fontenc}
\usepackage[utf8]{inputenc}
\usepackage[french]{babel}
\usepackage{lmodern}
\usepackage{amsmath,amssymb}
\usepackage[margin=2.35cm]{geometry}
\usepackage{booktabs,longtable,tabularx}
\usepackage{enumitem}
\usepackage{graphicx}
\usepackage{xcolor}
\usepackage{microtype}
\usepackage{float}
\usepackage{hyperref}

\hypersetup{colorlinks=true,linkcolor=blue!45!black,urlcolor=blue!45!black}
\setlist{nosep}
\newcommand{\code}[1]{\texttt{#1}}
\newcommand{\Kaut}{K^{*}_{\mathrm{aut}}}

\title{Prompt de recherche M4.3 (nom provisoire)\\
\large L'anti-corrélation entre queue des revenus d'intérêt et
criticité des avalanches est-elle structurelle, et quel mécanisme la
brise ?}
\author{Brouillon de travail --- successeur direct de M4.2B
(\code{m4\_2b\_credit\_soc})}
\date{6 août 2026 --- version de travail, non figée}

\begin{document}
\maketitle

\begin{quote}
\sloppy
\emph{Brouillon de travail, à discuter avant figeage. Successeur direct
de M4.2B (\code{m4\_2b\_credit\_soc/}), dont le prompt
(\code{prompts/PROMPT\_M4\_2B.md}) et le rapport final
(\code{report/rapport\_final.md}) sont des lectures préalables
obligatoires --- ce document ne répète pas ce qui y est déjà établi.}
\end{quote}

\tableofcontents

\section*{0. Ce que M4.2B a établi, et pourquoi ce prompt existe}
\addcontentsline{toc}{section}{0. Ce que M4.2B a établi, et pourquoi ce prompt existe}

M4.2B a testé si la queue supérieure de la distribution des intérêts
perçus pouvait être de type Pareto, de façon contrôlable, sans détruire
la structure en loi de puissance des avalanches de faillites. Le
programme n'a tranché ni oui ni non sur l'existence de la queue (limite
de puissance statistique à l'échelle testée, pas une réfutation) mais a
établi trois faits solides, confirmés sur graines disjointes :

\begin{enumerate}
\item \textbf{Épaisseur de queue et criticité des avalanches
  s'anti-corrèlent dans tout ce qui a été exploré.} $\eta$ aux deux
  extrêmes du sweep ($\rho=0{,}125$ et $\rho=4$) donne un $\hat\alpha$
  \emph{plus élevé} qu'à baseline (queue plus fine) alors que le
  rapport de branchement varie dans des sens opposés ($0{,}478$ vs
  $0{,}891$) ; le sweep $\delta=\sigma$ conjoint donne la queue la plus
  épaisse de toute la campagne ($\hat\alpha=3{,}05$) exactement là où
  le rapport de branchement s'effondre ($0{,}786\to0{,}386$). Aucune
  cellule des 37 testées ne combine queue plus épaisse ET criticité
  préservée ou accrue.
\item \textbf{La cible de principal arithmétique est structurellement
  anti-multiplicative} : elle homogénéise le capital
  (Gini($K$)$\approx0{,}05$) et la dispersion du revenu d'intérêt est
  dominée à $77\,\%$ par le nombre de contrats accumulés avec l'âge
  (\code{deg\_out}, un compteur), $23\,\%$ seulement par la
  taille/taux des contrats. Aucune ablation de mécanisme n'a été
  tentée dans M4.2B (question 16 du prompt original) --- c'est le
  plus gros degré de liberté resté inexploité.
\item \textbf{$\lambda$ (naissances, \code{Config.lam}) contrôle déjà
  la taille de population indépendamment de $K_0$} --- vérifié
  directement dans \code{m4\_2b/model.py}. La campagne de confirmation
  a fixé $\lambda=30$ partout et utilisé $K_0$ comme proxy de taille,
  ce que le rapport lui-même signale comme confondu (Gini $\times 8$,
  plancher de renouvellement $0\to0{,}8$, part \code{deg\_out} jusqu'à
  $0{,}24$ d'écart en même temps que $K_0$ varie). C'est un défaut de
  plan de campagne, pas une limite du moteur.
\end{enumerate}

Ce prompt ne redemande pas « existe-t-il une queue Pareto
contrôlable ». Il pose une question plus étroite et décidable :
\textbf{l'anti-corrélation observée entre l'épaisseur apparente de la
queue des intérêts et la criticité des avalanches est-elle une
propriété structurelle de cette classe de modèle (marché à pool
$k\equiv2$, faillite cancel+destroy, service d'intérêt perpétuel), ou
existe-t-il un paramètre ou un mécanisme qui la brise --- qui déplace
les deux quantités indépendamment, ou dans un sens conjointement
favorable ?}

\section*{1. Trois sous-questions, individuellement tranchables}
\addcontentsline{toc}{section}{1. Trois sous-questions, individuellement tranchables}

\begin{itemize}
\item \textbf{D1 --- cartographie.} Sur l'espace déjà couvert par
  M4.2B ($\eta$, $\gamma$, $K_0$, $\delta$, $\sigma$,
  \code{target\_rule}) et sur ce que ce prompt ajoute (\S4),
  l'anti-corrélation ($\hat\alpha$, $b$) tient-elle partout, ou
  existe-t-il des poches où elle s'inverse ou disparaît ?
\item \textbf{D2 --- contrôle indépendant.} Existe-t-il un paramètre
  ou un mécanisme qui déplace $\hat\alpha$ (ou la forme de la queue,
  voir \S3) sans déplacer $b$/$\hat\tau$ dans le sens opposé --- ou
  qui déplace les deux dans un sens conjointement favorable (queue
  plus épaisse, criticité égale ou supérieure) ?
\item \textbf{D3 --- robustesse.} Si une telle région existe (D2),
  résiste-t-elle au changement de graine, de fenêtre temporelle, et à
  la taille du système testée proprement (\S5, via $\lambda$, pas via
  $K_0$) ?
\end{itemize}

Un résultat négatif sur D2 (l'anti-corrélation est structurelle, aucun
levier trouvé ne la brise) est un résultat scientifique valide et
publiable --- ce n'est pas un échec de ce programme. C'est même le
résultat le plus probable au vu de M4.2B ; ne pas s'obstiner à le nier.

\section*{2. Existence de la queue de Pareto : hypothèse de travail dès le départ, pas un critère bloquant}
\addcontentsline{toc}{section}{2. Existence de la queue de Pareto : hypothèse de travail dès le départ, pas un critère bloquant}

M4.2B a dépensé une part importante de son budget (trois révisions
successives d'un test d'existence, \S4 du rapport final) à essayer de
\emph{prouver} que la queue est une vraie loi de puissance, avant de
finalement traiter l'existence comme hypothèse. \textbf{Ce prompt
saute directement à ce point de départ.}

\begin{itemize}
\item L'existence d'une loi de puissance pure sur la queue des
  intérêts n'est \textbf{pas} le critère de succès. Le critère est la
  réponse aux sous-questions D1--D3 sur le plan (paramètre de queue,
  criticité).
\item Caractériser la queue avec la famille qui s'ajuste le mieux et
  de façon reproductible --- loi de puissance pure, tronquée, avec
  coude --- sans chercher à trancher laquelle est « la vraie »
  au-delà de ce que les données permettent. Si un exposant de loi de
  puissance tronquée (ou un paramètre de coude/coupure) caractérise
  mieux et plus stablement la réponse aux paramètres que $\hat\alpha$
  seul, l'utiliser et le dire.
\item Ne pas répéter le test d'auto-similarité à seuil $\times 2$ de
  M4.2B tel quel (déjà démontré non calculable à cette échelle, \S4 du
  rapport final) --- si un test d'existence est utile quelque part
  dans ce programme, le motiver spécifiquement et vérifier d'abord
  qu'il est calculable avec le volume de données attendu
  ($n_{\mathrm{tail}}$ estimé avant de lancer, pas après coup).
\item \textbf{La statistique de queue utilisée pour D1/D2/D3 (\S1)
  doit être choisie et gelée pendant la phase pilote, avant le début
  de la cartographie D1} --- pas ajustée en cours de campagne selon
  ce qui « a l'air le mieux ». Choix justifié sur la reproductibilité
  inter-graines (écart-type le plus faible à paramètres fixés), pas
  sur la valeur elle-même. Une fois gelée ($\hat\alpha$ pur, exposant
  tronqué, ou paramètre de coupure), le plan (\S1) est défini sur
  cette statistique jusqu'à la fin du programme --- changer de
  statistique en cours de route a coûté trois révisions de protocole
  à M4.2B (\S4 de son rapport final), ne pas répéter.
\end{itemize}

\section*{3. Durée et taille de run : adaptatives, pas un budget fixé à l'avance}
\addcontentsline{toc}{section}{3. Durée et taille de run : adaptatives, pas un budget fixé à l'avance}

Pas de $T$ ou de $N$ cible fixé dans ce prompt. Chaque configuration
doit tourner jusqu'à ce que son propre temps caractéristique soit
mesuré et la fenêtre d'analyse positionnée en conséquence --- en
généralisant dès le départ ce que M4.2B a dû faire en rattrapage
(régression FOPDT sur le renouvellement du décile supérieur,
\code{scripts/renewal\_relaxation\_all\_runs.py}) :

\begin{itemize}
\item Étendre l'estimation du temps de relaxation à la variable
  d'intérêt central (queue des intérêts, pas seulement le
  renouvellement du décile supérieur utilisé comme proxy en M4.2B ---
  vérifier si son temps d'autocorrélation propre diffère, cf. limite
  notée au \S8 du rapport final de M4.2B).
\item \textbf{Checkpointing obligatoire dès le premier run} (pas de
  rattrapage possible sans lui, leçon déjà actée en mémoire de session
  après M4.2B) : sauvegarder l'état complet picklable de
  \code{Simulation} à intervalle régulier et en fin de run, pour
  permettre l'extension d'un run existant plutôt que sa relance
  complète si son temps de relaxation s'avère plus long que prévu.
\item Une cellule dont le temps de relaxation ne peut pas être atteint
  dans un budget de calcul raisonnable est documentée comme telle
  (sévère, non stationnaire confirmée) plutôt que forcée à un $T$
  arbitraire --- convention déjà en place en M4.2B, à garder.
\item Profiler le coût réel (s/pas, en fonction de $\lambda$ au
  minimum) avant tout engagement de campagne à plusieurs cellules ---
  pas après.
\end{itemize}

\section*{4. Où chercher un levier (D2) : ce qui est déjà disponible sans nouveau code}
\addcontentsline{toc}{section}{4. Où chercher un levier (D2) : ce qui est déjà disponible sans nouveau code}

Pas de liste de mécanismes pré-autorisée ni interdite --- la liberté
de diagnostic de M4.2B (\S17 de son prompt) est reconduite. Deux
précisions cependant :

\begin{itemize}
\item \textbf{Le contrôle géométrique
  (\code{target\_rule="geometric"}) est déjà implémenté, testé et
  parité-vérifiée avec M4.2 --- c'est un point de départ à coût quasi
  nul pour tester si le canal multiplicatif (que la cible arithmétique
  supprime, \S0.2) fait une différence sur le plan ($\hat\alpha$,
  $b$), avant d'écrire une seule ligne de mécanisme nouveau. Ce doit
  être le tout premier run de la phase pilote.} Le rapport M4.2B
  donne déjà un indice ($\hat\alpha=3{,}318$ en géométrique contre
  $3{,}890$ en arithmétique à baseline autrement identique, mais
  cellule sévère, un seul instantané) --- à re-mesurer proprement avec
  la méthode de fenêtrage adaptative (\S3) avant d'en tirer une
  conclusion.
\item \textbf{La décision d'ablation ne doit pas disparaître par
  défaut de temps.} M4.2B avait explicitement autorisé une refonte de
  mécanisme « en second recours » (\S17 de son prompt) et n'en a
  tenté aucune --- pas par choix motivé, mais parce que le temps a
  manqué. Ce prompt impose donc une règle de processus, pas une liste
  de mécanismes : \textbf{à un jalon fixé à l'avance (par exemple, la
  fin de la cartographie D1), une décision explicite et écrite doit
  être prise --- ablation nécessaire ou non, avec la justification ---
  avant de passer à la suite.} Si la réponse est « non nécessaire »,
  elle doit être aussi défendable et documentée qu'un « oui ».
\end{itemize}

\section*{5. Objectif B rigoureux : réutiliser la méthodologie déjà validée sur M4B, avec une réserve vérifiée}
\addcontentsline{toc}{section}{5. Objectif B rigoureux : réutiliser la méthodologie déjà validée sur M4B, avec une réserve vérifiée}

M4.2B n'a jamais testé l'indépendance de taille de l'exposant
d'avalanche correctement ($\lambda=30$ fixe, $K_0$ confondu comme
proxy de taille, \S8 de son rapport final). \code{recherche/sensibilite\_m4b/}
(campagne M4B, 594 runs, outils \code{lib\_metrics.py}/\code{lib\_screening.py}
déjà copiés sans changement dans \code{m4\_2b\_credit\_soc/scripts/}) a
résolu cette question \textbf{pour $k=3$} (\code{rapport\_final.tex}
\S« Avalanches~: loi tronquée\ldots » et table \code{tab:lois},
vérifié directement dans le source \LaTeX{}, pas de mémoire) :

\begin{itemize}
\item loi de puissance \textbf{tronquée} partout (LRT $p<10^{-6}$, pas
  de loi pure) ;
\item coupure $\hat s_c \propto \lambda^{1{,}5}$ (mesuré
  $\lambda=10\to30$ : $10{,}0\to51{,}6$ ; à $\lambda=100$ la coupure
  sort de la fenêtre observable, Vuong tronquée-vs-pure bascule
  franchement en faveur de la pure, $z=+21{,}0\pm0{,}9$, 5/5 graines) ;
\item exposant tronqué monotone avec $\lambda$
  ($1{,}584$~/~$2{,}041$~/~$2{,}313$), extrapolant à
  $\alpha_\infty\approx2{,}31$--$2{,}32$ --- la non-monotonie du fit
  en loi \emph{pure} est un artefact de coupure ;
\item $\lambda$ confirmé comme \textbf{pure taille finie} (intensivité
  $\pm2\,\%$ sur les observables normalisés).
\end{itemize}

\textbf{Réserve importante, à ne pas ignorer} : cette échelle en
$\lambda$ n'a été mesurée qu'à $k=3$. Le même rapport signale
explicitement, dans sa section « non observé/non confirmé », que
\textbf{la coupure $\hat s_c$ à $k=2$ (le \code{pool\_size} de M4.2B)
n'est pas confirmée en amplitude} (contraste exploratoire $-3{,}2$ vs
confirmé $+6{,}0\pm2{,}8$, métrique la plus bruitée du programme,
variance inter-graines $17\,\%$) et recommande explicitement de ne
piloter aucune conclusion sur la coupure ajustée seule --- de
préférer $b$ et $\hat\alpha$ (déclaré avec la taille de système) et
la susceptibilité, plus stables. Le lien entre M4B ($k=3$) et
M4.2B/ce prompt ($k\equiv2$) est donc une \textbf{hypothèse à
vérifier}, pas un fait transférable tel quel : la pente
$\lambda^{1{,}5}$ peut différer à $k=2$, même si le mécanisme
d'avalanche est hérité sans changement.

\textbf{Ce que ce prompt exige concrètement} :

\begin{itemize}
\item Réutiliser la méthodologie (loi tronquée par MLE, Vuong
  tronquée-vs-pure, échelle de $\lambda$ incluant au moins
  $\{10, 30, 100\}$ et un point au-delà si le budget le permet) ---
  pas les nombres de M4B.
\item \textbf{Ré-établir $\hat s_c(\lambda)$ sur le moteur M4.2B
  ($k\equiv2$)} avant de supposer une pente ; si elle diverge de
  $1{,}5$, le documenter, ne pas forcer l'accord.
\item Suivre la recommandation de M4B sur les observables de
  pilotage : \textbf{$b$ (rapport de branchement) et l'exposant
  tronqué $\hat\alpha$ déclaré avec sa taille de système sont les
  métriques primaires} du plan D1/D2/D3 côté avalanches ; la coupure
  ajustée reste une analyse secondaire/de confirmation, jamais le
  critère qui tranche à elle seule.
\item Toujours à paramètres par ailleurs identiques, jamais en
  faisant varier $K_0$ comme proxy de taille.
\end{itemize}

C'est un gain d'ingénierie direct (l'outillage existe déjà) et
méthodologique (fini les $\hat\tau$ indépendants par taille qui
avaient produit une dérive non interprétable dans le pilote M4.2B) ---
sans hériter une conclusion numérique qui n'a pas encore été vérifiée
dans ce régime.

\section*{6. Discipline scientifique (reconduite de M4.2B, sans changement)}
\addcontentsline{toc}{section}{6. Discipline scientifique (reconduite de M4.2B, sans changement)}

\begin{itemize}
\item Distinguer explicitement fait observé / inférence / hypothèse /
  incertitude à chaque affirmation quantitative.
\item Ne jamais sélectionner après coup graines, snapshots ou seuils
  donnant le résultat attendu.
\item Protocole figé après la phase pilote ; résultats négatifs
  conservés et rapportés comme tels (D1/D2 négatifs inclus, \S1).
\item Unité de réplication = snapshot pour la queue, run pour la
  moyenne de cellule, cellule pour l'effet de paramètre --- ne jamais
  fusionner ces échelles (leçon déjà actée : écart-type
  intra-instantané $\approx3{,}7\times$ l'écart-type inter-graines
  dans M4.2B, ne pas répéter la confusion).
\end{itemize}

\section*{7. Efficacité de calcul}
\addcontentsline{toc}{section}{7. Efficacité de calcul}

\begin{itemize}
\item Checkpointing dès le premier run (\S3) --- non négociable.
\item Profiler le coût par pas en fonction de $\lambda$ avant toute
  campagne à plusieurs cellules (\S3) ; documenter la mesure, pas une
  estimation.
\item Budget de parallélisme $\le6$ processus, un run indépendant par
  processus, conventions \code{cell\_id}/reprise déjà en place dans
  \code{scripts/lib\_lab.py} et équivalents (M4.2/M4.2B) ---
  réutiliser, ne pas reconstruire.
\end{itemize}

\section*{8. Parité et réutilisation d'outils}
\addcontentsline{toc}{section}{8. Parité et réutilisation d'outils}

\begin{itemize}
\item Moteur : partir de \code{m4\_2b/model.py} tel quel
  (l'institution arithmétique reste la référence sauf si D2/\S4
  justifie explicitement un changement, avec ablation propre et
  comparaison à la baseline arithmétique).
\item Figures et snapshots : réutiliser sans modification l'adaptateur
  \code{simulation\_lab} et le code de génération de figures hérité
  de M4.2/M4.2B (28 figures déjà validées) ; ajouter uniquement les
  figures propres au plan ($\hat\alpha$, $b$) et à la loi d'échelle
  $\hat s_c(\lambda)$.
\item Statistiques d'avalanches : réutiliser
  \code{lib\_metrics.py}/\code{lib\_screening.py} (\S5) sans les
  récrire.
\end{itemize}

\section*{9. Livrables}
\addcontentsline{toc}{section}{9. Livrables}

\begin{itemize}
\item Le plan ($\hat\alpha$ ou paramètre de queue équivalent,
  $b$/$\hat\tau$) cartographié sur l'espace testé (D1), avec figure
  dédiée (un point par cellule, graines agrégées, barres d'erreur
  inter-graines).
\item Verdict D2 explicite : levier trouvé ou non, avec le niveau de
  preuve exact (reproductibilité inter-graines, robustesse
  temporelle).
\item Si D2 positif : verdict D3 (robustesse de taille, méthodologie
  \S5).
\item Décision d'ablation (\S4) documentée par écrit, prise ou non
  prise, avec justification, à la date du jalon fixé.
\item Rapport autonome (contexte, grandeurs et méthode rappelés avant
  usage, comme le rapport final de M4.2B) + journal chronologique.
\end{itemize}

\section*{10. Question ouverte avant de figer ce prompt}
\addcontentsline{toc}{section}{10. Question ouverte avant de figer ce prompt}

Nom et emplacement de dossier (\code{m4\_3\_credit\_soc/} par défaut,
en suivant la convention M4$\to$M4B$\to$M4.2$\to$M4.2B --- à confirmer
ou changer), et tout ajustement de fond sur les sections ci-dessus.

\end{document}
